\chapter{Gestione delle eccezioni e dei file in Java}
\label{chap:eccezioni-file-java}

Il capitolo approfondisce il passaggio dei parametri e
l'organizzazione delle classi in package. Introduce inoltre le
eccezioni, utilizzate per segnalare e gestire condizioni anomale,
e le operazioni di lettura e scrittura dei file di testo, che
consentono di conservare i dati tra esecuzioni successive.%
\footnote{Materiale di riferimento:
S. Canazza, \emph{Java: gestione delle eccezioni e lettura/scrittura
di file}, \nolinkurl{2026_FI_11_canazza_Java_GestioneEccezioni_File.pdf}.}

\section{Passaggio dei parametri per valore}
\label{sec:fi-java-file-parametri}

\begin{definition}[Passaggio per valore]
\label{def:fi-java-file-passaggio-valore}
In Java, ogni argomento viene passato \emph{per valore}: il parametro
formale riceve una copia del valore dell'argomento. Assegnare un nuovo
valore al parametro non modifica la variabile del chiamante.
\end{definition}

Per i tipi primitivi si copia il dato; per i tipi riferimento si copia
il \emph{riferimento}, non l'oggetto
(\cref{sec:fi-java-riferimenti}). In quest'ultimo caso i due
riferimenti consentono di modificare lo stesso oggetto.

\begin{unipdcode}[language=Java]
public static void cambiaNumero(int n) {
    n = 20;
}

public static void aggiungiInteressi(BankAccount conto) {
    conto.deposit(100);
}
\end{unipdcode}

Dopo \texttt{cambiaNumero(x)}, il valore di \texttt{x} non cambia.
Dopo \texttt{aggiungiInteressi(account)}, invece, cambia il saldo
dell'oggetto indicato da \texttt{account}. Riassegnare un nuovo
\texttt{BankAccount} al parametro \texttt{conto} non cambierebbe
il riferimento del chiamante: il metodo modifica lo stato
condiviso dell'oggetto, non la variabile che contiene il riferimento.

\begin{figure}[htbp]
  \centering
  \begin{tikzpicture}[unipd/diagram,x=1cm,y=1cm]
  \node[unipd/node,align=center,minimum width=2.6cm]
    (v) at (-3.3,1.15) {Chiamante\\\texttt{x = 10}};
  \node[unipd/node,align=center,minimum width=2.6cm]
    (p) at (0.2,1.15) {Parametro\\\texttt{n = 10}};
  \draw[unipd/directed edge] (v.east) -- node[above,font=\scriptsize]{copia} (p.west);
  \node[unipd/node,align=center,minimum width=2.6cm]
    (r) at (-3.3,-0.8) {Chiamante\\\texttt{account: rif.}};
  \node[unipd/node,align=center,minimum width=2.6cm]
    (q) at (0.2,-0.8) {Parametro\\\texttt{conto: rif.}};
  \node[unipd/node,align=center,minimum width=2.45cm]
    (o) at (4.1,-0.8) {\texttt{BankAccount}\\\texttt{balance = 100}};
  \draw[unipd/directed edge] (r.east) -- node[above,font=\scriptsize]{copia} (q.west);
  \draw[unipd/directed edge] (q.east) -- (o.west);
  \draw[unipd/directed edge,densely dashed] (r.south) -- ++(0,-0.6)
    -| (o.south);
\end{tikzpicture}

  \caption{Il parametro riceve una copia del valore; una copia di un
  riferimento continua a individuare lo stesso oggetto.}
  \label{fig:fi-java-file-passaggio-valore}
\end{figure}

\clearpage

\begin{example}[Trasferimento fra oggetti]
Per trasferire un importo fra due conti, il metodo deve ricevere
il conto destinatario, non soltanto il suo saldo. Si aggiunge
alla classe \texttt{BankAccount} il seguente metodo:

\begin{unipdcode}[language=Java]
public void transfer(double amount, BankAccount other) {
    if (other == null || !(amount > 0) || !Double.isFinite(amount)
            || amount > balance) {
        throw new IllegalArgumentException();
    }
    withdraw(amount);
    other.deposit(amount);
}
\end{unipdcode}


\begin{unipdcode}[language=Java]
BankAccount corrente = new BankAccount(1000);
BankAccount risparmio = new BankAccount(200);
corrente.transfer(150, risparmio);
System.out.println(risparmio.getBalance()); // 350.0
\end{unipdcode}

\texttt{other.deposit(amount)} opera sull'oggetto destinatario.
Il metodo utilizza \texttt{withdraw()} e \texttt{deposit()}
introdotti nella \cref{sec:fi-java-classi-incapsulamento} e segnala
con \texttt{IllegalArgumentException} un importo non valido,
un saldo insufficiente o un destinatario nullo. Il lancio delle
eccezioni è approfondito nella \cref{sec:fi-java-file-eccezioni}.
\end{example}

\subsection{Invocazione di metodi di istanza}
\label{sec:fi-java-file-metodi-ricevente}

Un metodo di istanza richiede un \emph{ricevente}, cioè l'oggetto
su cui opera. In un metodo statico non esiste il riferimento implicito
\texttt{this}; per invocare un metodo di istanza occorre indicare
un oggetto (\cref{sec:fi-java-cicli-metodi-statici}).

\begin{unipdcode}[language=Java]
public class Contatore {
    private int valore;

    public void incrementa() {
        valore++;
    }

    public void incrementaDueVolte() {
        incrementa();       // equivale a this.incrementa()
        this.incrementa();
    }

    public static void main(String[] args) {
        Contatore c = new Contatore();
        c.incrementaDueVolte();
        System.out.println(c.valore);  // 2
        // incrementa(); // errore: manca il ricevente
    }
}
\end{unipdcode}

All'interno di \texttt{incrementaDueVolte()}, \texttt{this}
identifica l'oggetto \texttt{c} e può essere omesso quando
non sussistono ambiguità.

\subsection{Creazione e sostituzione di oggetti}
\label{sec:fi-java-file-ricostruire}

Un costruttore non è un metodo invocabile su un oggetto esistente:
\texttt{account.BankAccount()} non è valido. Per assegnare alla
variabile un nuovo oggetto si esegue nuovamente \texttt{new}:

\begin{unipdcode}[language=Java]
BankAccount account = new BankAccount();
BankAccount altroRiferimento = account;
account.deposit(500);
account = new BankAccount();
System.out.println(altroRiferimento.getBalance()); // 500.0
\end{unipdcode}

L'assegnamento cambia soltanto il riferimento di \texttt{account}:
l'oggetto precedente resta accessibile tramite \texttt{altroRiferimento}.

\section{I package di Java}
\label{sec:fi-java-file-package}

\begin{definition}[Package]
\label{def:fi-java-file-package}
Un \emph{package} è uno spazio dei nomi che raggruppa tipi Java
correlati. Permette di organizzarli e di distinguerli anche quando
hanno nomi semplici uguali.
\end{definition}

La dichiarazione \texttt{package} precede gli \texttt{import} e le
dichiarazioni dei tipi. Per convenzione, i nomi sono minuscoli e
spesso derivano da un dominio invertito:

\begin{unipdcode}[language=Java]
package com.example.conti;

public class BankAccount {
    // Definizione della classe
}
\end{unipdcode}

Il nome completo è \texttt{com.example.conti.BankAccount}. Di
norma, il file si trova nella directory relativa
\texttt{src/com/example/conti/BankAccount.java}; il corpo della
classe è omesso e si assume l'implementazione della
\cref{sec:fi-java-classi-incapsulamento}.
Un sottopackage come \texttt{com.example.conti.test} è distinto.

\subsection{Importazione e accessibilità}
\label{sec:fi-java-file-package-struttura}

La dichiarazione \texttt{import} consente di usare un tipo senza
ripeterne ogni volta il nome completo:

\begin{unipdcode}[language=Java]
package com.example.app;
import com.example.conti.BankAccount;

public class Main {
    public static void main(String[] args) {
        BankAccount conto = new BankAccount();
        System.out.println(conto.getBalance());
    }
}
\end{unipdcode}

\texttt{import com.example.conti.*;} rende disponibili i nomi dei
tipi accessibili del package, ma non importa i suoi sottopackage.
\texttt{java.lang} è importato implicitamente. L'importazione non
modifica l'accessibilità: senza modificatore, un membro è visibile
soltanto nello stesso package (\emph{package-private});
\texttt{protected} aggiunge specifiche possibilità di accesso per
le sottoclassi.

In assenza di dichiarazione \texttt{package}, il tipo appartiene al
\emph{package senza nome} (\emph{default package}). È adatto a
esempi isolati, non a progetti con più package: i tipi del package
senza nome non sono importabili da package con nome.

\subsection{Percorsi e compilazione}
\label{sec:fi-java-file-package-compilazione}

I percorsi possono essere \emph{assoluti} (dalla radice del
filesystem) oppure \emph{relativi} (dalla directory corrente).
L'opzione \texttt{-d} di \texttt{javac} indica la directory in cui
collocare i file \texttt{.class}, conservando l'organizzazione
prevista dai package:

\begin{unipdterminal}
$ mkdir -p bin
$ javac -d bin src/com/example/conti/BankAccount.java src/com/example/app/Main.java
$ java -cp bin com.example.app.Main
\end{unipdterminal}

L'esempio usa comandi di una shell Unix; su Windows può cambiare
la creazione delle directory. \texttt{-cp bin} indica la radice
delle classi compilate, mentre il nome dopo \texttt{java} è
qualificato con il package.

\section{Gestione delle eccezioni}
\label{sec:fi-java-file-eccezioni}

\subsection{Precondizioni e asserzioni}
\label{sec:fi-java-file-precondizioni}

\begin{definition}[Precondizione]
\label{def:fi-java-file-precondizione}
Una \emph{precondizione} è una condizione che deve valere prima
di un'operazione affinché ne sia garantito il comportamento
specificato.
\end{definition}

Il compilatore verifica i tipi degli argomenti, non vincoli come
la positività di un importo. Una precondizione può essere
specificata e verificata a esecuzione: la sua violazione può
essere segnalata con un'eccezione. Per esempio, una versione
controllata del metodo \texttt{BankAccount.deposit()} è:

\begin{unipdcode}[language=Java]
public void deposit(double amount) {
    if (!(amount > 0) || !Double.isFinite(amount)) {
        throw new IllegalArgumentException("Importo non valido");
    }
    balance += amount;
}
\end{unipdcode}

\texttt{System.exit(1)} termina il processo senza consentire
al chiamante di recuperare: per convenzione, \texttt{0} indica
successo e un codice non nullo un errore. L'istruzione
\texttt{assert condizione;} verifica invece un'assunzione interna;
la forma \texttt{assert condizione : messaggio;} aggiunge un
messaggio diagnostico. Le asserzioni sono disabilitate per impostazione
predefinita (abilitabili con \texttt{java -ea}): non devono
contenere operazioni con effetti collaterali e non sostituiscono
i controlli su input esterni.

\subsection{Lancio e gestione: \texttt{throw}, \texttt{try}, \texttt{catch}}
\label{sec:fi-java-file-throw-try-catch}

\begin{definition}[Eccezione]
\label{def:fi-java-file-eccezione}
Un'\emph{eccezione} è un oggetto appartenente alla gerarchia
di \texttt{Throwable}, utilizzato per segnalare una condizione
anomala. Le classi \texttt{Exception} ed \texttt{Error} ne
costituiscono due rami principali. L'istruzione
\texttt{throw} interrompe il normale flusso di esecuzione e
trasferisce il controllo a un gestore compatibile, eventualmente
in un metodo chiamante.
\end{definition}

\texttt{throw new IllegalArgumentException();} genera un'eccezione.
Se non viene intercettata localmente, si \emph{propaga} lungo la
catena delle chiamate.

\begin{definition}[Costrutto \texttt{try}/\texttt{catch}]
\label{def:fi-java-file-try-catch}
Il blocco \texttt{try} delimita le istruzioni protette; un blocco
\texttt{catch} gestisce un'eccezione il cui tipo è compatibile
con quello dichiarato. Dopo la cattura, l'esecuzione non riprende
nel punto in cui è stata lanciata l'eccezione.
\end{definition}

Data una stringa \texttt{riga}, la conversione può essere
gestita nel seguente modo:

\begin{unipdcode}[language=Java]
try {
    int n = Integer.parseInt(riga);
    System.out.println(n);
} catch (NumberFormatException e) {
    System.err.println("Intero non valido");
}
\end{unipdcode}

I blocchi \texttt{catch} possono essere più di uno: il primo tipo
compatibile viene selezionato, quindi quelli più specifici vanno
prima dei più generali. Se nessun gestore intercetta l'eccezione,
questa si propaga; nel thread principale può determinare una
terminazione anomala.

\begin{example}[Richiesta di un intero]
Il programma richiede un numero intero e ripete la lettura
finché non riceve un valore valido o termina il flusso di input.

\begin{unipdcode}[language=Java]
import java.util.Scanner;

public class LetturaIntero {
    public static void main(String[] args) {
        Scanner in = new Scanner(System.in);
        boolean valido = false;
        int numero = 0;
        while (!valido) {
            System.out.print("Inserisci un intero: ");
            if (!in.hasNextLine()) {
                break;
            }
            String riga = in.nextLine();
            try {
                numero = Integer.parseInt(riga);
                valido = true;
            } catch (NumberFormatException e) {
                System.err.println("Valore non valido.");
            }
        }
        if (valido) {
            System.out.println("Hai inserito: " + numero);
        }
    }
}
\end{unipdcode}


Una conversione non valida genera \texttt{NumberFormatException}
e consente di chiedere nuovamente il valore. Il controllo
\texttt{hasNextLine()} riconosce la fine del flusso.
\end{example}

L'oggetto eccezione espone, tra gli altri, \texttt{getMessage()}
e \texttt{printStackTrace()}. Quest'ultimo stampa normalmente
sullo \emph{standard error} la \emph{stack trace}, ossia la
sequenza delle chiamate che hanno condotto al lancio.
Una regola pratica è \emph{lanciare presto, catturare tardi}:
intercettare soltanto dove è possibile gestire il problema,
senza nasconderlo con blocchi \texttt{catch} vuoti.

\subsection{Eccezioni controllate e non controllate}
\label{sec:fi-java-file-checked}

Le classi delle eccezioni appartengono alla gerarchia di
\texttt{Throwable}, da cui derivano \texttt{Error} ed
\texttt{Exception}. Quest'ultima comprende, tra le altre,
\texttt{RuntimeException} e \texttt{IOException}:

\begin{figure}[htbp]
  \centering
  \begin{tikzpicture}[unipd/diagram,x=1cm,y=1cm]
  \tikzset{fi/exception/.style={unipd/node,align=center,
    minimum width=2.95cm,minimum height=0.64cm,
    inner xsep=4pt,inner ysep=5pt,font=\scriptsize}}
  \node[fi/exception] (thr)  at (0,4.75) {\texttt{Throwable}};
  \node[fi/exception] (err)  at (-3.875,3.45) {\texttt{Error}};
  \node[fi/exception] (exc)  at ( 3.875,3.45) {\texttt{Exception}};
  \node[fi/exception] (stack) at (-5.8,2.15) {\texttt{StackOverflowError}};
  \node[fi/exception] (memory)at (-1.95,2.15) {\texttt{OutOfMemoryError}};
  \node[fi/exception] (io)   at ( 1.95,2.15) {\texttt{IOException}};
  \node[fi/exception] (run)  at ( 5.8,2.15) {\texttt{RuntimeException}};
  \node[fi/exception] (file) at ( 1.95,0.85) {\texttt{FileNotFoundException}};
  \node[fi/exception] (arg)  at ( 5.8,0.85) {\texttt{IllegalArgumentException}};
  \node[fi/exception] (num)  at ( 5.8,-0.45) {\texttt{NumberFormatException}};
  \draw[unipd/directed edge] (err.north) -- (thr.south west);
  \draw[unipd/directed edge] (exc.north) -- (thr.south east);
  \draw[unipd/directed edge] (stack.north) -- (err.south west);
  \draw[unipd/directed edge] (memory.north) -- (err.south east);
  \draw[unipd/directed edge] (io.north) -- (exc.south west);
  \draw[unipd/directed edge] (run.north) -- (exc.south east);
  \draw[unipd/directed edge] (file.north) -- (io.south);
  \draw[unipd/directed edge] (arg.north) -- (run.south);
  \draw[unipd/directed edge] (num.north) -- (arg.south);
\end{tikzpicture}

  \caption{Gerarchia parziale delle classi derivate da
  \texttt{Throwable}; i rami \texttt{Error} e
  \texttt{RuntimeException} non sono controllati.}
  \label{fig:fi-java-file-gerarchia-eccezioni}
\end{figure}

\begin{definition}[Eccezioni checked e unchecked]
\label{def:fi-java-file-checked}
Le eccezioni \emph{controllate} (\emph{checked}) appartengono
a \texttt{Throwable} o a sue sottoclassi che non derivano da
\texttt{RuntimeException} né da \texttt{Error}. Il compilatore
impone di intercettarle mediante \texttt{catch} oppure di
dichiararle con \texttt{throws}.

Le eccezioni \emph{non controllate} (\emph{unchecked})
appartengono invece a \texttt{RuntimeException}, a \texttt{Error}
o alle rispettive sottoclassi: il compilatore non impone
di intercettarle o dichiararle.
\end{definition}

Le eccezioni controllate rappresentano generalmente condizioni
di errore che il programma può essere chiamato a gestire,
come problemi di accesso ai file. Quelle non controllate
comprendono errori di programmazione e condizioni anomale
dell'ambiente di esecuzione.

\texttt{IOException} è controllata; \texttt{IllegalArgumentException}
e \texttt{NumberFormatException} non lo sono.
La dichiarazione \texttt{throws} consente di propagare
un'eccezione al chiamante senza intercettarla.
La classificazione determina quindi un obbligo di
compilazione, non la gravità del problema.

\clearpage

\subsection{Eccezioni di I/O e propagazione con \texttt{throws}}
\label{sec:fi-java-file-io-throws}

L'apertura di un file può lanciare \texttt{FileNotFoundException},
sottoclasse di \texttt{IOException}. Un metodo che può produrre
un'eccezione controllata deve gestirla o dichiararla:

\begin{unipdcode}[language=Java]
import java.io.FileReader;
import java.io.IOException;
import java.nio.charset.StandardCharsets;

public class PrimoCarattere {
    public static void main(String[] args) throws IOException {
        try (FileReader in = new FileReader(
                "input.txt", StandardCharsets.UTF_8)) {
            int valore = in.read();
            if (valore != -1) {
                System.out.println((char) valore);
            }
        }
    }
}
\end{unipdcode}

\texttt{read()} restituisce il valore numerico di un'unità
\texttt{char} oppure \texttt{-1} a fine flusso: per questo il
programma verifica il risultato prima di convertirlo in
\texttt{char}. \texttt{throws IOException} dichiara una possibile
propagazione, mentre \texttt{throw} esegue un lancio. È possibile dichiarare
più tipi con \texttt{throws}, separandoli mediante virgole.
Usare \texttt{throws} anche in \texttt{main()} è lecito,
ma non gestisce l'errore.
\texttt{Scanner} appartiene a \texttt{java.util} e i suoi metodi
possono lanciare eccezioni non controllate, come
\texttt{InputMismatchException}.

\subsection{La clausola \texttt{finally}}
\label{sec:fi-java-file-finally}

Il blocco \texttt{finally} viene eseguito normalmente sia dopo
un \texttt{try} completato, sia prima della propagazione di
un'eccezione. È utile per il rilascio delle risorse:

\begin{unipdcode}[language=Java]
FileReader in = new FileReader("input.txt");
try {
    System.out.println(in.read());
} finally {
    in.close();
}
\end{unipdcode}

L'esempio richiede la gestione o dichiarazione di
\texttt{IOException}. Un'eccezione prodotta dalla chiusura può
mascherare quella precedente; per i file si preferisce quindi
\texttt{try-with-resources} (\cref{sec:fi-java-file-try-resources}).
\texttt{finally} non è garantito in caso di terminazione forzata
della JVM.

\clearpage

\section{Lettura e scrittura di file di testo}
\label{sec:fi-java-file-testo}

I file conservano dati oltre la terminazione del programma.
Un file di testo memorizza byte interpretati come caratteri
secondo una codifica, per esempio UTF-8
(\cref{sec:fi-caratteri}).

\subsection{Le API \texttt{java.io} e \texttt{java.nio}}
\label{sec:fi-java-file-io-nio}

\texttt{java.io} fornisce le classi tradizionali per i flussi di
byte e caratteri. \texttt{java.nio} e \texttt{java.nio.file}
offrono buffer, canali e operazioni moderne su percorsi e file.
Le API sono complementari; qui si utilizzano soprattutto i
lettori e gli scrittori di \texttt{java.io}.

\subsection{La classe \texttt{File}}
\label{sec:fi-java-file-classe-file}

Un oggetto \texttt{File} rappresenta un \emph{percorso}, non il
contenuto di un file aperto. Il diagramma indica i principali
metodi pubblici; \texttt{+} è la notazione UML per \emph{public}.

\begin{figure}[htbp]
  \centering
  \begin{tikzpicture}[unipd/diagram]
  \node[unipd/node,align=left,inner sep=7pt] {
    \begin{tabular}{@{}l@{}}
      \multicolumn{1}{c}{\textbf{\texttt{File}}}\\ \hline
      \texttt{+ exists(): boolean}\\
      \texttt{+ isFile(): boolean}\\
      \texttt{+ isDirectory(): boolean}\\
      \texttt{+ length(): long}\\
      \texttt{+ createNewFile(): boolean}\\
      \texttt{+ delete(): boolean}\\
      \texttt{+ renameTo(File): boolean}
    \end{tabular}
  };
\end{tikzpicture}

  \caption{Principali metodi di \texttt{java.io.File};
  notazione UML semplificata.}
  \label{fig:fi-java-file-uml-file}
\end{figure}

\begin{unipdcode}[language=Java]
File file = new File("dati.txt");
if (file.isFile()) {
    System.out.println(file.length()); // dimensione in byte
}
\end{unipdcode}

\texttt{createNewFile()} può lanciare \texttt{IOException};
\texttt{delete()} restituisce un \texttt{boolean} che indica
l'esito. I percorsi relativi si risolvono dalla directory di
lavoro del programma, non dalla posizione del sorgente.

\subsection{\texttt{FileReader}, \texttt{FileWriter} e buffer}
\label{sec:fi-java-file-classi-flussi}

\texttt{FileReader} legge caratteri, \texttt{FileWriter} li
scrive. \texttt{BufferedReader} e \texttt{BufferedWriter}
avvolgono tali flussi e aggiungono un buffer; il primo offre
\texttt{readLine()}, il secondo \texttt{newLine()}.
\texttt{PrintWriter} fornisce \texttt{print()}, \texttt{println()}
e \texttt{printf()} per scrivere testo formattato.

\clearpage

\begin{figure}[htbp]
  \centering
  \begin{tikzpicture}[unipd/diagram,x=1cm,y=1cm]
  \node[unipd/node,align=center,text width=4.1cm,inner sep=5pt]
    (fr) at (-2.95,2.2) {\textbf{\texttt{FileReader}}\\[-1pt]
    \rule{3.7cm}{0.25pt}\\[-2pt]\texttt{read(): int}};
  \node[unipd/node,align=center,text width=4.1cm,inner sep=5pt]
    (fw) at (2.95,2.2) {\textbf{\texttt{FileWriter}}\\[-1pt]
    \rule{3.7cm}{0.25pt}\\[-2pt]\texttt{write(String): void}};
  \node[unipd/node,align=center,text width=4.1cm,inner sep=5pt]
    (br) at (-2.95,0) {\textbf{\texttt{BufferedReader}}\\[-1pt]
    \rule{3.7cm}{0.25pt}\\[-2pt]\texttt{readLine(): String}};
  \node[unipd/node,align=center,text width=4.1cm,inner sep=5pt]
    (bw) at (2.95,0) {\textbf{\texttt{BufferedWriter}}\\[-1pt]
    \rule{3.7cm}{0.25pt}\\[-2pt]\texttt{newLine(): void}};
  \node[unipd/node,align=center,text width=4.1cm,inner sep=5pt]
    (pw) at (2.95,-2.75) {\textbf{\texttt{PrintWriter}}\\[-1pt]
    \rule{3.7cm}{0.25pt}\\[-2pt]\texttt{print(), println(), printf()}};
  \draw[unipd/directed edge] (br.north) -- node[right,font=\scriptsize]{usa} (fr.south);
  \draw[unipd/directed edge] (bw.north) -- node[right,font=\scriptsize]{usa} (fw.south);
  \draw[unipd/directed edge] (pw.north) -- node[right,font=\scriptsize]{può usare} (bw.south);
\end{tikzpicture}

  \caption{Metodi essenziali e composizione dei flussi di testo.
  Le frecce indicano l'utilizzo di un flusso sottostante,
  non l'ereditarietà.}
  \label{fig:fi-java-file-uml-flussi}
\end{figure}

\texttt{read()} restituisce un \texttt{int} oppure \(-1\)
a fine flusso. \texttt{readLine()} restituisce una riga senza
il terminatore oppure \texttt{null} a fine file; una riga vuota
produce invece \texttt{""}. Il buffering riduce gli accessi
al flusso sottostante.

\subsection{Lettura di un file}
\label{sec:fi-java-file-lettura}

Il programma legge \texttt{input.txt} riga per riga,
specificando la codifica UTF-8:

\begin{unipdcode}[language=Java]
import java.io.BufferedReader;
import java.io.FileReader;
import java.io.IOException;
import java.nio.charset.StandardCharsets;

public class LetturaFile {
    public static void main(String[] args) {
        try (BufferedReader in = new BufferedReader(
                new FileReader("input.txt", StandardCharsets.UTF_8))) {
            String riga;
            while ((riga = in.readLine()) != null) {
                System.out.println(riga);
            }
        } catch (IOException e) {
            System.err.println("Lettura fallita: " + e.getMessage());
        }
    }
}
\end{unipdcode}


In alternativa si può collegare \texttt{Scanner} a un
\texttt{FileReader} e usare \texttt{hasNextLine()} e
\texttt{nextLine()} (\cref{sec:fi-java-scanner}).

\subsection{Scrittura di un file}
\label{sec:fi-java-file-scrittura}

Il programma crea \texttt{risultati.txt}, oppure ne
sovrascrive il contenuto, usando un buffer:

\clearpage

\begin{unipdcode}[language=Java]
import java.io.BufferedWriter;
import java.io.FileWriter;
import java.io.IOException;
import java.nio.charset.StandardCharsets;

public class ScritturaFile {
    public static void main(String[] args) {
        try (BufferedWriter out = new BufferedWriter(
                new FileWriter("risultati.txt", StandardCharsets.UTF_8))) {
            out.write("Primo risultato");
            out.newLine();
            out.write("Secondo risultato");
            out.newLine();
        } catch (IOException e) {
            System.err.println("Scrittura fallita: " + e.getMessage());
        }
    }
}
\end{unipdcode}


\texttt{write()} non aggiunge una nuova riga; per questo si
invoca \texttt{newLine()}.

\begin{example}[Numerazione delle righe]
\label{ex:fi-java-file-righe-numerate}
Il seguente programma combina \texttt{BufferedReader} e
\texttt{PrintWriter} per numerare le righe di un file:

\begin{unipdcode}[language=Java]
import java.io.BufferedReader;
import java.io.FileReader;
import java.io.IOException;
import java.io.PrintWriter;
import java.nio.charset.StandardCharsets;

public class RigheNumerate {
    public static void main(String[] args) {
        try (BufferedReader in = new BufferedReader(
                new FileReader("input.txt", StandardCharsets.UTF_8));
             PrintWriter out = new PrintWriter(
                "numerate.txt", StandardCharsets.UTF_8)) {
            String riga;
            int numero = 1;
            while ((riga = in.readLine()) != null) {
                out.printf("%d: %s%n", numero++, riga);
            }
            if (out.checkError()) {
                System.err.println("Errore nella scrittura.");
            }
        } catch (IOException e) {
            System.err.println("Errore di I/O: " + e.getMessage());
        }
    }
}
\end{unipdcode}


Se \texttt{input.txt} contiene le righe \texttt{alfa} e
\texttt{beta}, l'output in \texttt{numerate.txt} è:
\begin{unipdterminal}
1: alfa
2: beta
\end{unipdterminal}
\end{example}

\begin{warning}[Errori di \texttt{PrintWriter}]
\texttt{PrintWriter} non propaga normalmente le
\texttt{IOException} di scrittura: \texttt{checkError()} verifica
l'eventuale stato d'errore. Anche \texttt{close()} non dichiara
\texttt{IOException}.
\end{warning}

\subsection{Chiusura automatica con \texttt{try-with-resources}}
\label{sec:fi-java-file-try-resources}

Il costrutto \texttt{try-with-resources} chiude automaticamente
le risorse dichiarate nell'intestazione che implementano
l'interfaccia \texttt{AutoCloseable}, anche in presenza di eccezioni:

\begin{unipdcode}[language=Java]
try (FileReader in = new FileReader("input.txt")) {
    System.out.println(in.read());
} catch (IOException e) {
    System.err.println(e.getMessage());
}
\end{unipdcode}

L'esempio riguarda la chiusura automatica e lascia implicita la
codifica; nelle operazioni effettive è opportuno specificarla.
Più risorse, separate da \texttt{;}, sono chiuse nell'ordine
inverso rispetto alla dichiarazione. Se la chiusura fallisce
mentre è già presente un'altra eccezione, quella di chiusura
viene registrata come \emph{eccezione soppressa}.

\subsection{Sovrascrittura, aggiunta e codifica}
\label{sec:fi-java-file-append}

\texttt{FileWriter} sovrascrive per impostazione predefinita.
Passando \texttt{true} nel costruttore, il contenuto viene invece
\emph{aggiunto} (\emph{append}):

\begin{unipdcode}[language=Java]
try (FileWriter out = new FileWriter(
        "log.txt", StandardCharsets.UTF_8, true)) {
    out.write("Nuova registrazione\n");
}
\end{unipdcode}

\texttt{StandardCharsets.UTF\_8} (da importare mediante
\texttt{java.nio.charset.StandardCharsets}) fissa la codifica.
L'operazione richiede la gestione o la dichiarazione di
\texttt{IOException}. In lettura e scrittura occorre usare
codifiche coerenti.

\subsection{Operazioni su file e directory}
\label{sec:fi-java-file-operazioni-avanzate}

Alcune operazioni sono disponibili tramite \texttt{File}; altre
tramite la classe \texttt{Files} di \texttt{java.nio.file}:

\begin{table}[htbp]
  \centering
  \small
  \begin{tabular}{@{}p{0.43\linewidth}p{0.49\linewidth}@{}}
    \toprule
    Metodo & Operazione \\
    \midrule
    \texttt{File.exists()} & Verifica che il percorso esista. \\
    \texttt{File.isDirectory()} & Verifica che sia una directory. \\
    \texttt{File.length()} & Dimensione del file in byte. \\
    \texttt{File.createNewFile()} & Crea un file se assente. \\
    \texttt{File.delete()} & Tenta di eliminare un file. \\
    \texttt{File.renameTo(File)} & Tenta di rinominare o spostare. \\
    \texttt{Files.copy(Path, Path, ...)} & Copia un file. \\
    \texttt{Files.move(Path, Path, ...)} & Sposta o rinomina un file. \\
    \bottomrule
  \end{tabular}
  \caption{Operazioni rappresentative sul filesystem.}
  \label{tab:fi-java-file-operazioni}
\end{table}

Le operazioni possono fallire per mancanza di permessi o
modifiche concorrenti al filesystem. Un controllo preventivo
con \texttt{exists()} non elimina la necessità di gestire
l'errore dell'operazione effettiva.

\subsection{Fine dell'input e tokenizzazione}
\label{sec:fi-java-file-eof-token}

L'acquisizione può terminare tramite una \emph{sentinella}
(un dato convenzionale, come \texttt{"Q"}) oppure raggiungendo
la \emph{fine del flusso} (EOF). Con \texttt{Scanner},
\texttt{hasNextLine()} restituisce \texttt{false} in quest'ultimo
caso. Da console, l'EOF è generalmente segnalato da
\texttt{Ctrl+D} sui terminali Unix e da \texttt{Ctrl+Z} seguito
da Invio su Windows.

Un \emph{token} è una parte del testo separata da delimitatori.
\texttt{Scanner} usa normalmente gli spazi bianchi; il metodo
\texttt{useDelimiter()} permette di impostare un'espressione
regolare differente.

\begin{unipdcode}[language=Java]
import java.util.Scanner;

public class ContaToken {
    public static void main(String[] args) {
        Scanner in = new Scanner(System.in);
        int conteggio = 0;
        while (in.hasNextLine()) {
            String riga = in.nextLine();
            try (Scanner parole = new Scanner(riga)) {
                while (parole.hasNext()) {
                    parole.next();
                    conteggio++;
                }
            }
        }
        System.out.println("Token: " + conteggio);
    }
}
\end{unipdcode}


Per separare, per esempio, le parti di
\texttt{"Anna/Rossi/18"}, si può usare:

\begin{unipdcode}[language=Java]
Scanner parti = new Scanner("Anna/Rossi/18");
parti.useDelimiter("/");
\end{unipdcode}

\clearpage

\subsection{Reindirizzamento, pipe e flussi standard}
\label{sec:fi-java-file-redirezione}

Java espone tre flussi standard: \texttt{System.in} (ingresso),
\texttt{System.out} (risultati) e \texttt{System.err}
(diagnostica). Gli ultimi due sono di tipo \texttt{PrintStream}.
La shell può collegarli a file senza modificare il programma:

\begin{unipdcode}[language=Java]
import java.util.Scanner;

public class Pappagallo {
    public static void main(String[] args) {
        Scanner in = new Scanner(System.in);
        while (in.hasNext()) {
            System.out.println(in.next());
        }
    }
}
\end{unipdcode}


\begin{unipdterminal}
$ javac Pappagallo.java
$ java Pappagallo < testo.txt > parole.txt
$ java Pappagallo < testo.txt > parole.txt 2> errori.txt
\end{unipdterminal}

\texttt{<} reindirizza l'ingresso, \texttt{>} l'uscita e
\texttt{2>} il flusso di errore. \texttt{>} crea o sovrascrive
il file, mentre \verb|>>| aggiunge in fondo. Una \emph{pipe}
(\texttt{|})
collega l'uscita standard di un programma all'ingresso di un altro:

\begin{unipdterminal}
$ java Pappagallo < testo.txt | sort > ordinato.txt
$ java Pappagallo < testo.txt | sort -fu | wc -l
\end{unipdterminal}

Il secondo esempio usa le utilità Unix \texttt{sort} e \texttt{wc}
per contare le righe distinte ignorando le maiuscole. La sintassi
precisa dipende dalla shell. Separare \texttt{System.err} da
\texttt{System.out} evita di includere i messaggi diagnostici
nei dati trasferiti attraverso una pipe.
